\chapter{Relations $r_{\theta_{k}}$}

	\noindent
	We start from the relationship Flja\_Abstract:\\\\
	$Flja\_Abstract: SNEIs \rightarrow LCF_{2p}$\\\\
	It is a set of R-pairs.
	\\\\
	
	\noindent
	As it is a set, we can do a partition from it. And the subsets we will obtain, will be relations themself too. They are subsets of R-pairs.
	\\\\
	
	\noindent
	Here is where we are going to understand the "r" in the columns of the tables of the previous chapter. THEY ARE RELATIONS. \textbf{Each column is an $R_{\theta_{k}}$ relation.}
	\\\\
	
	\noindent
	Flja\_Abstract relation associates, in each column, to each SNEI, $\aleph_{0}$ lcfs. And we we have $\aleph_{0}$ columns. We could call this THE MEGAPACK.
	\\\\
	
	\noindent
	The relation $r_{\theta_{k}}$. Is exactly the same algorithm as Flja\_Abstract, but using only the column K. All lcfs inside that column belong to the same universe $\theta_{k}$. A Pack with $\aleph_{0}$ different members. A SUBPACK in Flja\_Abstract, A PACK in $r_{\theta_{k}}$.
	\\\\
	
	\noindent
	Each relation:  $r_{\theta_{k}}: SNEIs \rightarrow \theta_{k}$\\\\
	Uses ALL SNEIs as Domain, but only the universe $\theta_{k}$ as the set where it obtains, all PACKs for the SNEIs, from. INSIDE that concrete $r_{\theta_{k}}$. All universes $\theta_{k}$ are disjoint between them, so the same R-pair is NOT in two $r_{\theta_{k}}$ at the same time. Defining Flja\_Abstract, we have defined, IN ONE STEP, all the infinite list of different relations $r_{\theta_{k}}$.
	\\\\
	
	\noindent
	\textless ** \textit{See unnecesary comment 0015, for some examples if needed} \textgreater
	\\\\